Jesus describes a man working in a field who must leave his coat behind.
The danger is so urgent that going back for something useful would put his
life at risk. A person on a roof must leave without gathering possessions
from the house. Believing Jesus means acting on his warning, even when
obedience costs something. He also speaks of pregnant and nursing mothers,
for whom escape will be especially difficult. Their distress is part of
what he tells his disciples to expect.

The disciples have asked about the Temple's destruction and about the
Lord's coming. His directions for flight concern Judea. His promise that
the angels will gather his people reaches to the ends of the heavens.
Saint John Chrysostom distinguishes these events. The local emergency
requires escape; the final coming brings a gathering that Christ himself
accomplishes.

Before that coming there will be people claiming to offer access to him.
They will announce that Christ is in a desert or a hidden room. They will
even offer impressive signs. Jesus tells his disciples beforehand to refuse
them. He compares his own coming to lightning visible across the sky.
There will be no need for an informant who knows where the returning Lord
is concealed.

Someone longing for deliverance could be drawn toward the person who
promises it most confidently. Jesus requires a judgment beyond that longing.
A claim must be tested against what he has said, even when it offers
something we deeply want. The disciples' desire for rescue needs his
direction. So does ours.

Later in this Mass, the Communion antiphon gives his assurance about
prayer: ``believe that you shall receive.'' We hear those words after a
Gospel in which people dear to God must endure great distress.
Confidence in prayer cannot mean that suffering proves a failure to
believe. Jesus says that the days are shortened for the elect's sake;
the elect are the people he has chosen. His care for them is real
while their danger is real.

At the beginning of Mass, we heard God's promise of peace to the
exiles in Babylon. Jeremiah tells those same people to build houses,
plant gardens and pray for the city where they are living. Their
return will not be immediate. They have work to do in the place
from which they hope to be delivered. The false prophets' promise
of a quick return would encourage them to neglect that work.
God's promise enables them to continue it. We can ask him for
deliverance while remaining responsible for the good he has
given us to do. Sometimes that good must be done through a
long wait, with the strength that has to be asked for again.

Paul describes a gift already received by the Colossians. The Father
has rescued them from darkness and brought them into his beloved
Son's kingdom. They have redemption through Christ's Blood and the
forgiveness of sins. Yet Paul keeps praying for them. He asks that
they know God's will, do good work and have strength to persevere.
Forgiveness has opened a life in which they still need to grow.

He gives thanks to the Father who has ``made us worthy'' to share
the saints' inheritance. Saint Thomas Aquinas explains why Paul
says made. The worthiness itself is something God has brought about
in them. They have not first made themselves fit and then persuaded
him to give them a place. Their very capacity to receive his gift
belongs among the things for which they thank him.

That makes Paul's prayer for good works hopeful. God can act in
people who still need to learn obedience. Knowledge of his will is
to become conduct pleasing to him. The strength Paul asks for
includes patience when doing good is difficult and perseverance
when the promised good is delayed. These people need more than
the assurance that they once received something. They need God's
continued help to live faithfully.

In the Collect, we ask God to arouse the wills of his faithful.
We ask for his action within people who already belong to him.
Our willing and doing matter; the prayer asks that they bear
fruit. Admitting that our wills need his help is therefore a
truthful way to pray as believers. The Father who has made us
worthy can enable the obedience we still find difficult.

When we bring our offerings, the Secret asks God to accept them
with our prayers and turn our hearts toward him. It asks for a
change from earthly cravings to heavenly desires. Bringing a
gift does not mean that everything we want is already right.
We are asking him to change what we seek as well as to hear
what we ask.

Aquinas explains prayer as the mind's ascent to God. In asking
another person for something, we may have to persuade him to
become willing. God's goodness already precedes our asking.
Prayer raises us toward him. He really gives what Paul asks
for: knowledge, strength and endurance. Receiving those gifts
changes what a person is able to understand, want and do.

The Communion's promise belongs to Jesus' teaching in Mark
about faith in God. Just after that promise, Jesus tells
those standing in prayer to forgive anyone against whom
they have a grievance. He speaks the assurance and the
command to the same disciples. Their treatment of another
person belongs to their response to God.

Saint Bede explains this command through faith working
by love. He also addresses Christians who have not reached
extraordinary perfection. They can ask for pardon with
confidence as they forgive those who have sinned against
them. You do not need the power to work wonders before
you can ask for mercy. You do need to surrender the
refusal to show mercy to someone else.

This is a demanding answer to the assurance that we
shall receive. A petitioner seeking pardon may want
God's mercy while wanting to keep a grievance untouched.
Jesus requires forgiveness within the act of praying.
The change we need can reach the very desire with
which we approached him.

If you are holding a grievance, ask God for that
person's good, with a willingness to receive his
help in forgiving. The wrong done to you need not
determine everything you will ask for that person.
The Father's action in making his people worthy
gives you reason to seek help even where your
own will resists.

At Communion we receive Christ, through whose
Blood we have the redemption Paul proclaims.
The prayer that follows reception asks that
the mysteries we have received cure what
remains faulty in us. We have received a
gift whose healing we still need. Remaining
faults can be acknowledged and submitted to
his cure without calling his generosity
into question.

The inheritance for which Paul gives thanks
is life with God among his saints. Christ
promises to gather his people, and he
prepares them to live with him. His
forgiveness and his help in changing
your heart are already gifts of that
salvation. You can ask for the further
healing you need with confidence in
the Lord who intends to bring you
into its fullness.
